\documentclass[letterpaper]{article}

\usepackage[letterpaper,margin=1in]{geometry}

% Compile from the repository root so this path resolves to the pinned WSP
% submodule.
\input{wsp/testing/test-case-library.tex}

\begin{document}
\def\TCID{TC-0004}
\def\TCTitle{Windows console host}
\def\TCPurpose{Verify that the console host shows the cells of a screen
(characters through code page 437, attributes without the blink bit) and its
cursor in a console screen buffer.}
\def\TCPriority{\TCPriorityHigh}

\def\TCRequirementRef{REQ-0003}
\def\TCDesignRef{Not applicable.}
\def\TCFeatureRef{\texttt{include/wfc/brun71/console.hpp};
\texttt{src/console.cpp}; \texttt{tests/console\_test.cpp}.}
\def\TCTestTechnique{Scenario test: draw on a screen, show it, read the
console buffer back.}

\def\TCPreconditions{The library is built for Windows and the test program
\texttt{brun71\_console\_test} exists.}
\def\TCEnvironment{Windows. The process needs a console; without one the
program exits 77 and CTest reports the test as skipped. The screen buffer is
never made active, so the console the test runs in is not disturbed. No key
input is read (it would take keys typed into that console). CTest runs the
program with a 60 second timeout.}
\def\TCAssumptions{Not applicable.}

\def\TCInputData{Box-drawing, letter and glyph characters in yellow on blue;
a blinking colour; the cursor hidden and shown with a full-height shape; a
40-column screen.}
\def\TCInitialState{A new screen buffer, not active.}

\def\TCProcedure{\begin{enumerate}
  \item Run \texttt{brun71\_console\_test} (CTest \texttt{TC-0004-console},
    in WFC \texttt{wfc-brun71-TC-0004-console}).
  \item Check the character and attribute shown at the printed cells and at
    a blank one, and the cursor position and visibility.
  \item Check that a blinking colour is shown without the blink bit and that
    the cursor shows with its full-height shape.
  \item Check a cell in the last column of a 40-column screen.
\end{enumerate}}
\def\TCExpectedResult{The program exits with status 0 and prints
\texttt{console: ok} (or exits 77 without a console); no \texttt{FAIL} line is
written.}
\def\TCPostConditions{The screen buffer is closed.}

\TCTemplate

\end{document}
